\documentclass[10pt,a4paper]{amsart}
\usepackage{amsmath,amssymb,amsthm,graphicx}
% Retained minimal preamble for the Stacks Project chapter topology.tex (Topology).
% The Stacks Project uses one single shared preamble for every chapter; the snapshot stages it
% identically as sources/stacks-topology/preamble.tex, sources/stacks-chow/preamble.tex and
% sources/stacks-etale-cohomology/preamble.tex.  The block below is copied mechanically, line for
% line, from that shared preamble: the Xy-pic setup (lines 14-18) and the native theorem
% declarations together with the mathematical macros (lines 154-217), plus the project's own
% reference/slogan/history display environments (lines 8-12, via the verbatim package).
% Nothing was added, removed or redefined; the wrapper class and the amsmath/amssymb/amsthm
% base packages are supplied by the assembler.
% Chapter: topology.tex -- Topology.
\usepackage[all]{xy}

% We use 2cell for 2-commutative diagrams.
\xyoption{2cell}

% Theorem environments.
%
\theoremstyle{plain}
\newtheorem{theorem}[subsection]{Theorem}
\newtheorem{proposition}[subsection]{Proposition}
\newtheorem{lemma}[subsection]{Lemma}

\theoremstyle{definition}
\newtheorem{definition}[subsection]{Definition}
\newtheorem{example}[subsection]{Example}
\newtheorem{exercise}[subsection]{Exercise}
\newtheorem{situation}[subsection]{Situation}

\theoremstyle{remark}
\newtheorem{remark}[subsection]{Remark}
\newtheorem{remarks}[subsection]{Remarks}

\numberwithin{equation}{subsection}

% Macros
%
\def\lim{\mathop{\mathrm{lim}}\nolimits}
\def\colim{\mathop{\mathrm{colim}}\nolimits}
\def\Spec{\mathop{\mathrm{Spec}}}
\def\Hom{\mathop{\mathrm{Hom}}\nolimits}
\def\Ext{\mathop{\mathrm{Ext}}\nolimits}
\def\SheafHom{\mathop{\mathcal{H}\!\mathit{om}}\nolimits}
\def\SheafExt{\mathop{\mathcal{E}\!\mathit{xt}}\nolimits}
\def\Sch{\mathit{Sch}}
\def\Mor{\mathop{\mathrm{Mor}}\nolimits}
\def\Ob{\mathop{\mathrm{Ob}}\nolimits}
\def\Sh{\mathop{\mathit{Sh}}\nolimits}
\def\NL{\mathop{N\!L}\nolimits}
\def\CH{\mathop{\mathrm{CH}}\nolimits}
\def\proetale{{pro\text{-}\acute{e}tale}}
\def\etale{{\acute{e}tale}}
\def\QCoh{\mathit{QCoh}}
\def\Ker{\mathop{\mathrm{Ker}}}
\def\Im{\mathop{\mathrm{Im}}}
\def\Coker{\mathop{\mathrm{Coker}}}
\def\Coim{\mathop{\mathrm{Coim}}}

% Boxtimes
%
\DeclareMathSymbol{\boxtimes}{\mathbin}{AMSa}{"02}

%
% Macros for moduli stacks/spaces
%
\def\QCohstack{\mathcal{QC}\!\mathit{oh}}
\def\Cohstack{\mathcal{C}\!\mathit{oh}}
\def\Spacesstack{\mathcal{S}\!\mathit{paces}}
\def\Quotfunctor{\mathrm{Quot}}
\def\Hilbfunctor{\mathrm{Hilb}}
\def\Curvesstack{\mathcal{C}\!\mathit{urves}}
\def\Polarizedstack{\mathcal{P}\!\mathit{olarized}}
\def\Complexesstack{\mathcal{C}\!\mathit{omplexes}}
% \Pic is the operator that assigns to X its picard group, usage \Pic(X)
% \Picardstack_{X/B} denotes the Picard stack of X over B
% \Picardfunctor_{X/B} denotes the Picard functor of X over B
\def\Pic{\mathop{\mathrm{Pic}}\nolimits}
\def\Picardstack{\mathcal{P}\!\mathit{ic}}
\def\Picardfunctor{\mathrm{Pic}}
\def\Deformationcategory{\mathcal{D}\!\mathit{ef}}

% --- appended 2026-10-12 (Chief): the project own definitions of the reference/slogan/history
% display environments, copied mechanically from the shared Stacks preamble lines 8-12
% (\usepackage{verbatim} plus the three \newenvironment lines, which the chapter preamble keeps
% verbatim from the project preamble). Nothing added, removed or redefined: these environments
% carry no mathematical content and the project itself typesets them as comments. Without these
% three lines a counted statement whose own text uses them fails with Environment slogan undefined.
\usepackage{verbatim}
\newenvironment{reference}{\comment}{\endcomment}
\newenvironment{slogan}{\comment}{\endcomment}
\newenvironment{history}{\comment}{\endcomment}
\title{A quotient of a sober space by an equivalence relation with finite fibres and lifting generalizations is Kolmogorov}
\author{The Stacks Project Authors}
\date{stacks-project:topology.tex:local-snapshot-20261012; CC BY 4.0}
\begin{document}
\maketitle
\noindent\textit{Editorial note.} All delivered mathematical text is the author's own, copied from the pinned Stacks Project chapter source sources/stacks-topology/pinned-original-source.tex; no statement or proof step was rewritten, reordered or removed.\par\medskip
\noindent\textit{Editorial note.} Preamble provenance: the Stacks Project keeps one single shared preamble for every chapter, staged here as sources/stacks-topology/preamble.tex; this card uses retained-preamble.tex, which copies that shared preamble mechanically (the Xy-pic setup, the native theorem declarations, the mathematical macros and the reference/slogan/history display environments). Nothing was added, removed or redefined; the wrapper class and the base packages are supplied by the assembler.\par\medskip
\noindent\textit{Editorial note.} Bibliography provenance: the chapter carries no chapter-local bibliography (the project keeps one shared bibliography file that is not staged here) and neither the counted statement nor its proof contains any citation, so the card carries no bibliography and no entry was reconstructed.\par\medskip
\noindent\textit{Editorial note.} This card was selected from the staged candidate list for this chapter (candidates-halfA.json / candidates-halfB.json), where it is the candidate labelled lemma-quotient-kolmogorov; the statement and proof spans are the candidate's own statement\_lines and proof\_lines, checked against the source so that each span opens and closes a real, uncommented environment.\par\medskip
\noindent\textit{Allowed standard prerequisites (editorial).}\begin{itemize}
\item The Kolmogorov separation axiom (distinct points are separated by an open set one way or the other) and the specialization relation u leads to u-prime between points. Set up in the author's sections on specialization and on Kolmogorov spaces, before the counted statement. Fixes the conclusion of the statement and the relation used throughout the proof.
\item The hypotheses of the statement: openness of the map, sobriety of the source, finiteness of the fibres of the two maps, lifting of generalizations along them, and the quotient description of the target. The five numbered items of the counted statement itself (source lines 3936-3943). Each hypothesis is used in the proof: finiteness makes equivalence classes finite, lifting of generalizations transfers specializations, and openness relates the topology of the quotient to that of the source.
\end{itemize}\medskip
\begin{lemma}
\label{lemma-quotient-kolmogorov}
Suppose that $s, t : R \to U$ and $\pi : U \to X$ are continuous maps
of topological spaces such that
\begin{enumerate}
\item $\pi$ is open,
\item $U$ is sober,
\item $s, t$ have finite fibres,
\item generalizations lift along $s, t$,
\item $(t, s)(R) \subset U \times U$ is an equivalence relation on $U$ and
$X$ is the quotient of $U$ by this equivalence relation (as a set).
\end{enumerate}
Then $X$ is Kolmogorov.
\end{lemma}

\begin{proof}
Properties (3) and (5) imply that a point $x$ corresponds to an
finite equivalence class $\{u_1, \ldots, u_n\} \subset U$ of the equivalence
relation. Suppose that $x' \in X$ is a second point corresponding to
the equivalence class $\{u'_1, \ldots, u'_m\} \subset U$.
Suppose that $u_i \leadsto u'_j$ for some $i, j$. Then for any
$r' \in R$ with $s(r') = u'_j$ by (4) we can find $r \leadsto r'$
with $s(r) = u_i$. Hence $t(r) \leadsto t(r')$. Since
$\{u'_1, \ldots, u'_m\} = t(s^{-1}(\{u'_j\}))$ we conclude that
every element of $\{u'_1, \ldots, u'_m\}$ is the specialization of
an element of $\{u_1, \ldots, u_n\}$.
Thus $\overline{\{u_1\}} \cup \ldots \cup \overline{\{u_n\}}$ is
a union of equivalence classes, hence of the form $\pi^{-1}(Z)$
for some subset $Z \subset X$. By (1) we see that $Z$ is closed in $X$
and in fact $Z = \overline{\{x\}}$ because $\pi(\overline{\{u_i\}})
\subset \overline{\{x\}}$ for each $i$. In other words, $x \leadsto x'$
if and only if some lift of $x$ in $U$ specializes to some lift of
$x'$ in $U$, if and only if every lift of $x'$ in $U$ is a specialization
of some lift of $x$ in $U$.

\medskip\noindent
Suppose that both $x \leadsto x'$ and $x' \leadsto x$. Say $x$ corresponds
to $\{u_1, \ldots, u_n\}$ and $x'$ corresponds to $\{u'_1, \ldots, u'_m\}$
as above. Then, by the results of the preceding paragraph, we can find a
sequence
$$
\ldots \leadsto u'_{j_3} \leadsto u_{i_3} \leadsto u'_{j_2} \leadsto
u_{i_2} \leadsto u'_{j_1} \leadsto u_{i_1}
$$
which must repeat, hence by (2) we conclude that
$\{u_1, \ldots, u_n\} = \{u'_1, \ldots, u'_m\}$, i.e., $x = x'$.
Thus $X$ is Kolmogorov.
\end{proof}

\end{document}
